\chapter{Introduction}\label{ch:introduction}

\section{Motivation}

The quadratic time and memory complexity of self-attention
\cite{vaswani2017attention} has long defined the practical ceiling
of sequence modelling. As context windows expand toward millions
of tokens, this ceiling translates into prohibitive deployment
costs in memory, latency, and energy. The field has consequently
shifted toward sub-quadratic alternatives, and linear-time
recurrent models, in particular State Space Models (SSMs) and
the RWKV lineage, have moved from academic curiosities into a
production substrate for frontier-scale systems
\cite{li2025minimax01,liu2025hunyuanturbos}.

This industrial adoption is evidenced by several recent
deployments. MiniMax-01 demonstrated that linear attention can
drive a $456$\,B-parameter mixture-of-experts model with context
windows reaching $4$\,M tokens \cite{li2025minimax01}. Tencent's
Hunyuan-TurboS \cite{liu2025hunyuanturbos} and NVIDIA's
Nemotron-H family \cite{blakeman2025nemotronh} replace the
majority of self-attention layers with Mamba blocks, reporting up
to a threefold inference speed-up at parity. Further examples
include AI21's Jamba \cite{lieber2024jamba}, the Zamba2 suite
\cite{glorioso2024zamba2}, and the Evo~2 genomic foundation model,
which uses a linear convolutional multi-hybrid to reach a $1$\,M
nucleotide context \cite{brixi2026evo2}. Now that these models
form a production substrate, their structural limitations
translate directly into measurable deployment costs.

Underneath this architectural variety, the models share a
recurrent core and linear-time token mixing. They also share
characteristic failure modes: an inability to resolve local
structure at the correct scale, unreliable associative recall,
and restricted recurrent dynamics. Numerous fixes have been
proposed, including quad-directional spatial shifts in vision,
depthwise convolutions in audio, and decorrelation steps in
language modelling, yet they have not been compared on a single
substrate at matched compute, which makes their relative gains
difficult to attribute to architectural cause. This thesis
addresses the missing cross-cut by treating
\textit{mechanism-prior alignment} as a hypothesis: the gain a
mechanism delivers should be predictable from how much of the
targeted deficit the underlying backbone leaves uncovered.

\section{Contributions}

The thesis introduces three functional mechanisms that address
documented expressivity deficits of linear-time recurrent
models: a Multi-Scale Depthwise Convolution that addresses the
short-range temporal-hierarchy axis, a Damped Harmonic
Oscillator that extends the recurrent transition class beyond
real-diagonal, and a Chunked Value Decorrelation that targets
associative memory under interference. The mechanisms are
evaluated on a cross-architecture matrix (RWKV-6, Mamba-2,
Linear Attention) across two access modes (causal, LION
bidirectional), two parameter scales (7M, 14M), and three
probes: LibriSpeech for the principal axis-1 evaluation, Common
Voice English for cross-distribution validation within axis~1,
and the Multi-Query Associative Recall benchmark for axis-2
isolation at sequence lengths $T \in \{64, 256, 1024\}$.

The resulting evidence supports the central claim of the thesis:
a mechanism's gain on a task is proportional to the residual
structural deficit the architecture leaves uncovered along the
axis the mechanism targets, modulated by what the
task-as-exercised-by-data exercises. Both productive transfer
(deficit-proportional MSDC, BREAK-and-NULL DHO, asymmetric CVD)
and a single counter-example (LION-LIT~$\times$~CVD on Linear
Attention, where the absence of decay inverts the transfer sign)
populate the matrix and validate mechanism-prior alignment as a
predictive criterion. The strongest single empirical statement
of the thesis is the inversion of the standard
sub-quadratic-versus-quadratic narrative on Multi-Query
Associative Recall at $T = 1024$, where every linear-time
backbone augmented with MSDC or CVD solves the task while the
causal Transformer baseline fails.

\section{Research questions}\label{sec:intro:rqs}

The empirical evaluation answers four research questions, each
mapped to a specific section of the experimental chapter.

\paragraph{RQ1 (framework).} Can the expressivity of causal
linear-time recurrent architectures be characterised through an
axis decomposition under which a mechanism's gain on a task is
predictable from the residual structural deficit the
architecture leaves uncovered along the axis the mechanism
targets?

\paragraph{RQ2 (causal matrix).} Do the three proposed mechanisms
transfer across the three causal backbones in a way that orders
with the architecture-deficit map, and do the three exhibit
distinct transfer signatures rather than a single common
pattern?

\paragraph{RQ3 (bidirectional adaptation).} Does the LION
mechanism-level bidirectional wrapper preserve the transfer
pattern observed in the causal mode, and are there structural
prerequisites in the bidirectional setting that are absent in
the causal setting?

\paragraph{RQ4 (cross-task and cross-distribution).} Does the
same mechanism-prior alignment hold when the matrix is exercised
on a different acoustic distribution (Common Voice) and on a
synthetic axis-2 probe (Multi-Query Associative Recall), and
what modulations does the task-as-exercised-by-data introduce?

\section{Structure of the thesis}\label{sec:intro:structure}

The remainder of the thesis is organised as follows.
Chapter~\ref{ch:related} surveys the architectural and
mechanism literature on which the proposed mechanisms build, the
formal expressivity bounds that delimit the diagonal class, and
the evaluation probes used by the empirical chapter. The chapter
positions the thesis within the wider linear-time recurrent
literature.
Chapter~\ref{ch:theoretical} establishes the theoretical
foundation, introduces the axis decomposition that frames RQ1,
formalises the three causal linear-time backbones within a
unified recurrence template, presents the LION parallel form
together with the decay-class taxonomy, and presents the three
mathematical primitives on which the proposed mechanisms build.
Chapter~\ref{ch:proposed} introduces the three proposed
mechanisms, the unified LION wrapper, and the diagnostic probes
used at the parameter level.
Chapter~\ref{ch:experiments} reports the empirical matrix,
answers RQ2, RQ3, and RQ4 across LibriSpeech, Common Voice, and
Multi-Query Associative Recall, and synthesises the
mechanism-prior alignment criterion.
Chapter~\ref{ch:conclusions} summarises the thesis's
contributions and identifies the diagonal-to-diagonal-plus-low-rank
extension as the natural next step for the axis-3 question that
the present scope does not address.
